\chapter{Architecture générale}

\section{Vue d'ensemble}

L'architecture sépare les traitements hors ligne, réalisés lors de la préparation des connaissances, et le traitement en ligne, exécuté à chaque question.

\begin{figure}[H]
\centering
\resizebox{\textwidth}{!}{%
\begin{tikzpicture}[node distance=1.5cm and 1.2cm]
  \node[source] (xmi) {Métamodèle\\XMI};
  \node[traitement,right=of xmi] (json) {Conversion\\XMI vers JSON};
  \node[stockage,right=of json] (uml) {Graphe UML\\Neo4j};

  \node[source,below=of xmi] (docs) {Fiches\\textuelles};
  \node[traitement,right=of docs] (chunks) {Découpage et\\embeddings};
  \node[stockage,right=of chunks] (index) {Documents, chunks\\et index Neo4j};

  \node[traitement,below=of chunks] (extract) {Extraction et\\validation métier};
  \node[stockage,right=of extract] (metier) {Concepts métier\\validés};

  \node[source,right=2.0cm of uml] (question) {Question et\\historique};
  \node[traitement,below=of question] (retriever) {Retriever\\GraphRAG};
  \node[traitement,below=of retriever] (llm) {Mistral\\génération};
  \node[stockage,below=of llm] (answer) {Réponse\\sourcée};

  \draw[fleche] (xmi)--(json); \draw[fleche] (json)--(uml);
  \draw[fleche] (docs)--(chunks); \draw[fleche] (chunks)--(index);
  \draw[fleche] (docs.south)--(extract.west); \draw[fleche] (extract)--(metier);
  \draw[fleche] (question)--(retriever); \draw[fleche] (uml)--(retriever);
  \draw[fleche] (index)--(retriever); \draw[fleche] (metier)--(retriever);
  \draw[fleche] (retriever)--(llm); \draw[fleche] (llm)--(answer);
\end{tikzpicture}}
\caption{Architecture globale du projet}
\label{fig:architecture-globale}
\end{figure}

\section{Composants}

\begin{table}[H]
\centering
\caption{Responsabilités des composants principaux}
\begin{tabularx}{\textwidth}{>{\ttfamily}p{0.31\textwidth}X}
\toprule
Composant & Responsabilité \\
\midrule
XmiToJson.py & Convertir le métamodèle XMI en JSON exploitable. \\
GrapheBuilder\_v3.py & Construire le graphe UML, importer et indexer les documents. \\
graphrag\_retriever.py & Résoudre les entités et sélectionner sous-graphes et passages. \\
Agent\_LG.py & Construire le prompt final et appeler Mistral. \\
api.py & Authentification, conversations, messages et routes HTTP. \\
pipeline\_graphe\_metier.py & Proposer et contrôler les concepts et relations métier. \\
validation\_metier.py & Appliquer les contrôles déterministes et corrections bornées. \\
graphe\_metier\_builder.py & Importer uniquement une sélection explicitement validée. \\
\bottomrule
\end{tabularx}
\end{table}

\section{Choix d'une base unique à deux couches}

Le graphe technique et le graphe métier sont conservés dans la même base afin de permettre leur traversal conjoint. Ils restent cependant séparés par leurs labels et leurs relations : une \code{Classe} UML n'est jamais fusionnée avec un \code{ConceptMetier}. Le pont explicite \code{REPRESENTE\_PAR} maintient la distinction entre représentation technique et sens fonctionnel.

\capture{architecture-application.png}{Architecture réelle de l'application ou diagramme de composants}{fig:capture-architecture}

